%! AP Statistics — Unit 5, Lesson 5.1 — Guided Notes
\documentclass[10pt]{article}
\usepackage{apstats-article}
\usepackage{apstats-boxes}
\newcommand{\TallMath}[1]{$\displaystyle #1\rule[-1.4em]{0pt}{3.2em}$}

\begin{document}

\pageheader{Unit 5, Lesson 5.1}{Guided Notes}
\namedateperiod

\vspace{0.15in}

\begin{objectivebox}
By the end of this lesson, I will be able to:
\begin{enumerate}
  \item \writeline
  \item \writeline
\end{enumerate}
\end{objectivebox}

\vspace{0.1in}

\begin{vocabbox}
\termblanklong{Statistic}
\termblanklong{Parameter}
\termblanklong{Sampling variability}
\termblanklong{Sampling distribution}
\end{vocabbox}

\vspace{0.1in}

\begin{hookbox}
\textbf{Polling scenario:} Three agencies each survey 500 randomly selected voters
about a ballot measure and report support rates of 54\%, 51\%, and 56\%.

Does this mean one agency is wrong? \writeline

What is the most likely explanation for the differences? \writeline
\end{hookbox}

\vspace{0.1in}

\begin{notesbox}{Statistics vs.\ Parameters}
A \textbf{parameter} is \blank{8cm} --- it describes the \blank{3.5cm}.\\[6pt]
A \textbf{statistic} is \blank{8cm} --- it describes a \blank{3cm}.\\[10pt]

\textbf{Example:} A survey of 200 voters finds that 54\% support a measure.
\begin{itemize}
  \item The \textbf{statistic} is \blank{6cm}
  \item The \textbf{parameter} is \blank{6cm}
\end{itemize}
\end{notesbox}

\vspace{0.1in}

\begin{notesbox}{Sampling Variability}
When we take many samples of the same size from the same population, the statistics we
compute \blank{8cm} from sample to sample.

This is called \blank{6cm} and is \blank{3cm} (expected / unexpected).

\textbf{Key idea (VAR-1.G.1):} Variation in statistics for samples taken from the same
population may be \blank{8cm}.
\end{notesbox}

\vspace{0.1in}

\begin{notesbox}{The Sampling Distribution}
The \textbf{sampling distribution} of a statistic is the distribution of \blank{8cm}
\writeline
for all possible \blank{6cm} of a given size from a population.

\textbf{Three things to distinguish:}
\begin{enumerate}
  \item Distribution of the \textbf{population}: describes every individual value.
  \item Distribution of a \textbf{sample}: describes the values in one sample.
  \item \textbf{Sampling distribution} of a statistic: describes \writeline
\end{enumerate}

\textit{Think of it this way:} One sample $\to$ one dot. All possible samples $\to$
the whole \blank{5cm}.
\end{notesbox}

\vspace{0.1in}

\begin{practicebox}
\textbf{Quick Check.} In each situation, identify whether the value described is a
statistic or a parameter.

\begin{enumerate}
  \item The average age of all students at your school. \hfill \blank{3.5cm}
  \item The proportion of red cards in a random draw of 10 from a standard deck.
        \hfill \blank{3.5cm}
  \item The true mean body temperature of all healthy adults. \hfill \blank{3.5cm}
\end{enumerate}
\end{practicebox}

\end{document}
